\documentclass[11pt]{article}

% ── Packages ──────────────────────────────────────────────────────────────────
\usepackage[UTF8,heading=false,fontset=none]{ctex}
\setCJKmainfont{Microsoft YaHei}
\setCJKsansfont{Microsoft YaHei}
\linespread{1.08}
\usepackage[left=0.5in, right=0.5in, top=0.2in, bottom=0.2in]{geometry}
\usepackage{enumitem}
\usepackage{titlesec}
\usepackage{hyperref}
\usepackage{xcolor}

% ── Colours ───────────────────────────────────────────────────────────────────
\definecolor{navy}{RGB}{28, 57, 100}
\definecolor{rulecol}{RGB}{28, 57, 100}

% ── Hyperlink setup ───────────────────────────────────────────────────────────
\hypersetup{
  colorlinks=true,
  urlcolor=navy,
  linkcolor=black
}

% ── Page style ────────────────────────────────────────────────────────────────
\pagestyle{empty}
\raggedbottom
\raggedright
\setlength{\parindent}{0pt}
\setlength{\parskip}{0pt}

% ── Section heading ───────────────────────────────────────────────────────────
\titleformat{\section}{
  \vspace{1pt}\color{navy}\large\bfseries\sffamily
}{}{0em}{}[{\color{rulecol}\titlerule[0.6pt]}\vspace{1pt}]

\titlespacing*{\section}{0pt}{9pt}{4pt}

% ── Bullet list ───────────────────────────────────────────────────────────────
\newlist{rbullet}{itemize}{1}
\setlist[rbullet]{
  label=\textbullet,
  leftmargin=1.4em,
  itemsep=0pt,
  parsep=0pt,
  topsep=0pt,
  partopsep=0pt,
  after=\vspace{0pt}
}

% ── Job header helper ─────────────────────────────────────────────────────────
% \jobheader{Company}{Title}{Period}{Location}
\newcommand{\jobheader}[4]{%
  \vspace{1pt}%
  {\textbf{#1} \hfill \textit{#3}} \\[-2pt]
  {\textit{#2} \hfill \textit{#4}} \\[-4pt]
}

% ── Project header helper ─────────────────────────────────────────────────────
\newcommand{\projheader}[2]{%
  \vspace{1pt}%
  {\textbf{#1} \hfill \textit{#2}} \\[-4pt]
}

% ══════════════════════════════════════════════════════════════════════════════
\begin{document}

% ── Name ──────────────────────────────────────────────────────────────────────
\begin{center}
  {\fontsize{18}{20}\selectfont\bfseries\sffamily 陈雨荃}\\[1pt]
  \small
  158-2118-4124 \;|\;
  \href{mailto:chenyuquan297@hotmail.com}{chenyuquan297@hotmail.com} \;|\;
  \href{https://linkedin.com/in/yuquan-althea-chen}{linkedin}
\end{center}

% ══ 教育背景 ════════════════════════════════════════════════════════════════
\section{教育背景}

{\textbf{卡内基梅隆大学} \hfill 2025年8月 -- 2026年12月}\\[-2pt]
{\textit{信息系统管理硕士 (MISM)} \hfill 美国 匹兹堡}\\[1pt]
{\small 相关课程：数据分析Python (A)、颠覆性技术管理 (A$-$)、自然语言处理 (B+)、统计学 (A$-$)}\par

\vspace{1pt}

{\textbf{南京大学} \hfill 2021年9月 -- 2025年6月}\\[-2pt]
{\textit{信息管理与信息系统 学士} \hfill 中国 南京}\\[1pt]
{\small 绩点：4.32/5.0 | 相关课程：机器学习与深度学习实践 (94)、社会网络分析 (95.2)、文本数据管理 (93)、大数据系统 (92)、运筹学 (89)}\par

% ══ 工作经历 ══════════════════════════════════════════════════════════
\section{工作经历}

\jobheader{Whelix —— AI机器人初创公司（生物科技方向）}{AI工程师实习生}{2026年5月 -- 至今}{美国 波士顿}
\begin{rbullet}
  \item 从零设计并上线了一套生产级多智能体AI平台，直接基于原生\textbf{Gemini SDK}开发（不依赖现成的智能体框架），掌控每一步工具调用逻辑；整套平台共用一套后端服务，支撑\textbf{四个}与Slack深度集成的智能体，其中两个是仅限CEO使用、可查看实时销售管道与融资数据的专属助手，部署在\textbf{Google Cloud Run}
  \item 通过\textbf{22个自定义工具函数}将平台与\textbf{四个}外部系统（Slack、Notion、Confluence、Google Workspace）打通，实现结构化数据的读写；设计了一套所有智能体（顶层及各专属智能体）共用的工具路由层，在执行前将每个请求缩小至最相关的工具子集以提升工具调用准确率
  \item 构建了双层记忆架构——实时结构化事实写入，配合基于GCS的对话日志与夜间LLM驱动的信息整合——使无状态的webhook调用也能维持连贯的多轮对话上下文
  \item 通过定时Cloud Run触发器，端到端自动化了公司\textbf{三项}高频重复性工作流程，并独立编写了一套涵盖\textbf{11个数据库}的知识架构方案，将公司分散的内部文档整合为统一的信息记录系统
\end{rbullet}

\jobheader{巴斯夫聚氨酯特种产品（中国）有限公司}{数字化转型实习生，亚太区单体事业部}{2024年5月 -- 2025年2月}{中国上海}
\begin{rbullet}
  \item 发现巴斯夫MDI/TDI定价流程中的一项关键信息缺口：市场分析师需人工监测3--4个跨两种语言的新闻信息源以追踪需求信号，该流程易出错，却直接影响到巴斯夫在中国市场拥有较高定价话语权背景下的定价决策
  \item 设计并部署了一套\textbf{端到端、基于LLM（GPT-4o）的市场情报引擎}，实现双语新闻监测、信息提取与内容整合的全流程自动化——完全取代人工流程，并直接向定价决策者交付结构化的周报与月报
  \item 与市场部相关人员紧密合作，确保输出格式符合高层汇报要求，使该工具从上线之初即被采纳使用，彻底消除了团队人工新闻筛选的需求
  \item 支持\textbf{8个以上}其他数字化转型项目，涵盖供应链、环境健康安全（EHS）及生产领域——包括Power BI仪表盘重设计与数据管道自动化，在业务需求与落地技术方案之间承担技术转化角色
\end{rbullet}


% ══ 精选项目与获奖经历 ══════════════════════════════════════════════
\section{部分项目与获奖经历}

\projheader{本科毕业论文：基于决策导向学习的AUUC最大化因果推断方法研究}{2025年2月 -- 5月}
\begin{rbullet}
  \item 提出了一套\textbf{决策导向学习（DFL）框架}，通过复合损失函数直接优化营销ROI，并建模为整数线性规划（ILP）问题——在Criteo因果推断基准数据集上，相较基线方法实现\textbf{3\%}的增量效果提升，在个体处理效应估计上展现出更优的稳定性
\end{rbullet}

\projheader{CCF大数据与计算智能大赛 \textnormal{（1,557名参赛者）}}{2023年10月 -- 12月}
\begin{rbullet}
  \item 作为队长带领团队\textbf{获得全国前三（二等奖）}；负责团队任务分配，并在机器学习与运筹学两种技术路线间进行权衡决策，以优化库存决策方案
\end{rbullet}


% ══ 技术技能 ═════════════════════════════════════════════════════════
\section{技术技能}

{\small\noindent
\textbf{数据分析与人工智能：} LLM智能体系统 (LangGraph, AutoGen)、机器学习、增量建模 (Uplift Modeling)、因果推断、自然语言处理、时间序列预测\\[1pt]
\textbf{技术栈：} Python (Pandas, NumPy, Pytorch, Scikit-learn)、SQL、Power BI、KNIME、R\\[1pt]
\textbf{语言能力：} 英语（工作水平，雅思7.5）、普通话（母语）、西班牙语（初级）
}

\end{document}
